% Generated draft: requires mathematical and source review.
\documentclass[11pt]{article}
\usepackage[margin=25mm]{geometry}
\usepackage{fontspec}
\setmainfont{Latin Modern Roman}
\usepackage{amsmath,amssymb,mathtools,unicode-math}
\setmathfont{Latin Modern Math}
\usepackage{enumitem,needspace,xurl,hyperref}
\setlength{\parindent}{0pt}
\setlength{\parskip}{5pt}
\setlength{\emergencystretch}{3em}
\newcommand{\N}{\mathbb{N}}
\newcommand{\Z}{\mathbb{Z}}
\newcommand{\Q}{\mathbb{Q}}
\newcommand{\R}{\mathbb{R}}
\newcommand{\C}{\mathbb{C}}
\newcommand{\F}{\mathbb{F}}
\newcommand{\A}{\mathbb{A}}
\newcommand{\PP}{\mathbb{P}}
\newcommand{\E}{\mathbb{E}}
\newcommand{\eps}{\varepsilon}
\newcommand{\dd}{\,\mathrm{d}}
\begin{document}
\section*{UnsolvedMath ID: 700491}
\textbf{Source problem code:} catalogue entry
\subsection{Definitions and mathematical statement}
Let \(G\) be a finite group with identity \(1\). An involution is an element of order exactly \(2\). A normal set of involutions is a subset \(D\subseteq G\) consisting of involutions and invariant under conjugation: \(g^{-1}Dg=D\) for every \(g\in G\). For \(x\in G\), define
\[
C_D(x)=\{d\in D:dx=xd\}.
\]
Classify the finite groups \(G\) for which there exist such a set \(D\) and a subset \(\varnothing\ne T\subsetneq D\) satisfying
\[
\begin{aligned}
&\forall a\in T, &&C_D(a)\subseteq T,\\
&\forall a,b\in T, &&\bigl(ab=ba\ \text{and}\ ab\ne1\bigr)\Longrightarrow C_D(ab)\subseteq T.
\end{aligned}
\]
Here normality refers to conjugation invariance of \(D\); no assumption that \(D\) generates \(G\), that \(D\) is a single conjugacy class, or that \(T\) is conjugation invariant is supplied.
\subsection{Short English statement}
Determine which finite groups admit a conjugation-invariant set \(D\) of involutions with a nonempty proper subset \(T\) such that every member of \(D\) commuting with an element of \(T\) lies in \(T\), and every member of \(D\) commuting with the nonidentity product of two commuting elements of \(T\) also lies in \(T\).
\subsection{Sources}
\begin{itemize}
\item Local UnsolvedMath registry, permanent ID \(700491\), source code \texttt{catalogue entry}, title \emph{Kourovka Notebook Problem 8.41}. The supplied record labels the problem open; this status has not been independently verified.
\item \emph{The Kourovka Notebook (Unsolved Problems in Group Theory)}, Problem 8.41, proposed by A. A. Makhn\"ev in the 8th issue (1982). The supplied background identifies the 21st issue edition (2026), \texttt{arXiv:1401.0300v48}. The supplied collection description names E. I. Khukhro and V. D. Mazurov as editors. These bibliographic details are supplied attribution, not independently verified source consultation.
\end{itemize}
\end{document}
